%============================================================
\chapter{事前登録した含意の一覧}
\label{app:prereg}
%============================================================

本稿は理論から導いた含意を、観測にあたる前に文章として固定した。
支持・棄却の別を後から選べないようにするためである。
本文は結果のみを述べる。ここに登録の文言と判定を一覧する。

登録は22件、支持6件（うち部分的支持3件）、棄却11件、検証不能5件である。

\section{一覧}

\begin{center}
\begin{tabularx}{\textwidth}{@{}llXl@{}}
\toprule
記号 & 対象 & 登録した含意と結果 & 判定 \\
\midrule
A & 族2 & 乖離と取引頻度。水準は支持だが測れたのは $\kappa$ & 部分的支持 \\
B & 族5 & 返戻金と可測性。符号が一意に定まらない & 検証不能 \\
C & 手数料 & 外部化の対価。機能で説明できない反例も発見 & 部分的支持 \\
D & 無人化 & 依存の数と無人期間。$\Phi$1 と同一の含意であり同時に棄却 & 棄却 \\
E & 法人統計 & $\CCC$ の先行性（$p=0.72$） & 棄却 \\
F & 法人統計 & 規模と与信の向き。符号が逆 & 棄却 \\
J & 手数料 & 料率差は相関 $\rho$ に対応。$\rho$ が観測できない & 検証不能 \\
K & $\phi_{\rm cog}$ & $\delta^\ast>\delta$。四カテゴリで符号が負 & 棄却 \\
L & 容量制約 & 高齢層の契約選択。消費者側データなし & 検証不能 \\
M & 法人統計 & $\CCC$ 上昇の構成効果。業種内が83\% & 棄却 \\
N & 法人統計 & $\Delta\ln W$ と $\Delta\ln r$ の正相関（$\beta=0.83$） & 支持 \\
O & 法人統計 & 遅れ構造。年次データで検出できず & 棄却 \\
P & 法人統計 & ラグを含む総弾力性。$0.830\to0.821$ & 棄却 \\
T & 労働側 & 新法による $\kappa$ の移動。統計が存在しない & 検証不能 \\
U & 個人企業 & 営業利益率が役員報酬を戻した値に近い & 支持 \\
X & 族5 & 手数料率と離職率の負相関。307社で無相関 & 棄却 \\
Y & 族5 & 規模と離職率の相関の職種順序。全国では再現されず、定性的な区分のみ成立 & 部分的支持 \\
Z & 族5 & B1 の大きい職種ほど返戻金が強い。順位相関 $\rho=0.000$ & 棄却 \\
$\Phi$1 & 手続き & 依存の数と改修までの期間。複製の集中が結果を作っていた & 棄却 \\
$\Phi$2 & 手続き & $M(t)$ は途絶に先行して減少（$p=3.2\times10^{-9}$） & 支持 \\
$\Phi$3 & 手続き & $\CCC$ の非負性。企業部門も全業種で非負であり対照が立たない & 検証不能 \\
$\Phi$4 & 手続き & 族7 の型。稀である点は当たり、残る型が外れた & 棄却 \\
\bottomrule
\end{tabularx}
\captionof{table}{事前登録した22件の含意と判定。}
\label{tab:prereg-all}
\end{center}

G・H・I は制度による料率規制を対象としたが、
母集団が3件しかないことが判明したため設計段階で取り下げた。
Q・R・S は労働者視点の検討で登録したが、検証に着手していない。

\section{登録が判定を変えた箇所}

事前に固定していなければ後付けの解釈を採っていたと考えられる箇所を記録する。

\begin{itemize}
\item 族5において、上限制手数料が99.7\%消滅している事実は、
可測性仮説の傍証として解釈することが可能であった。
実際には上限制の価格水準が一桁低いだけで無関係である。

\item 第\ref{part:industry}章において、$\CCC$ と売上高の同期負相関を
「$\CCC$ は景気指標として機能する」と読み、
大企業の与信ポジションを「大企業が取引先を支えている」と解釈しえた。
いずれも後付けの物語である。

\item $\Phi$2 において、途絶した $\Phi$ の残高が事前に減少していることだけを見れば
支持と読めた。継続した $\Phi$ にも同じ窓を当てると、180 日では両者が分離しない。
対照を置くと決めていなければ、全体の減少基調を支持と読み違えた。

\item $\Phi$4 において、唯一現れた族7 が受益者の選定を手続きの外へ出していることは、
命題\ref{prop:beneficiary}の傍証として提示できた。
しかし登録した文言は「残るのは 7-4」であり、現れたのは 7-2 である。
機構が命題と整合することは、外れた含意を救わない。
\end{itemize}

\section{登録の運用で生じた不備}

Dと$\Phi$1 は同一の含意を二度登録したものである。
第13段階と第16段階で別々に登録し、その時点で重複に気づいていなかった。
登録の一覧を章ごとに持ち、全体で突き合わせていなかったことによる。

%============================================================
